% =====================================================================
%  The Patient Quant — back matter
%  Glossary, notes and references, appendices A-E, index,
%  acknowledgments, about the author, colophon.
% =====================================================================

\backmatter

% ---- Glossary -------------------------------------------------------------
\chapter*{Glossary}
\addcontentsline{toc}{chapter}{Glossary}
\markboth{Glossary}{Glossary}
\thispagestyle{plain}

\noindent\small
Terms are defined as they are used in this book. Where the field uses a term
in a narrower sense than ordinary English, the finance sense is given.
Symbols that recur throughout the text are given with their meaning.

\begin{description}[style=nextline,leftmargin=1.1em,labelindent=0pt,
  font=\normalfont\ttfamily\color{accent},itemsep=3pt,parsep=0pt]

\item[Arithmetic mean] The simple average of a series of returns. Always
larger than the geometric mean by approximately half the variance per
period.

\item[Certainty equivalent] The risk-adjusted value of an uncertain
expected return. Under quadratic utility it is the expected return less
the risk-aversion coefficient times half the variance of the estimate.

\item[Compound annual growth rate] The constant annual rate that would
produce the same terminal value as the actual series of returns. The
money a saver actually earns, as opposed to the arithmetic mean.

\item[Deflated Sharpe Ratio] A test of whether an observed Sharpe ratio
is larger than the expected maximum Sharpe ratio of a family of trials,
corrected for the non-normality of the returns.

\item[Drawdown] The fall of a portfolio's value from its running peak.
The book distinguishes the depth of the worst drawdown from the time it
took to recover.

\item[Ensemble] The set of model configurations run in parallel and
combined by a weighted average, where the weights depend on the members'
realised records.

\item[Forward return] The total return realised between a formation date
and a later date, at the model's horizon.

\item[Geometric mean] The compounded average of a series of returns. Equal
to the arithmetic mean less approximately half the variance per period.

\item[Kelly criterion] The position size that maximises the expected
logarithmic growth of wealth, given a known edge.

\item[Live window] The period over which a fund had a price on the panel
and could therefore have been held. A backtest that ignores the live
window has a survivorship bias.

\item[Normal inverse-gamma] A conjugate prior pair under which a linear
regression with an unknown scale has a closed-form posterior for both the
coefficients and the scale.

\item[Posterior variance of the mean] The variance of the \emph{estimate}
of the expected return, as opposed to the variance of the return itself.
The first governs how much to believe a forecast; the second governs how
much an asset can lose.

\item[Probability of backtest overfitting] The probability that the
configuration that wins on one half of a record loses on the other half.
A value above one half means the selection procedure is harmful.

\item[Reality check] A bootstrap test of whether the best member of a
family of strategies beats a benchmark by more than the family would by
chance.

\item[Sequence of returns] The order in which returns arrive. A saver's
outcome depends on it even though a buy-and-hold investor's does not.

\item[Sharpe ratio] Excess return divided by its standard deviation. The
book reports both the naive ratio and Lo's corrected ratio.

\item[Sleeve] One of the four asset groups --- United States equity,
international equity, bonds, real assets --- into which the strategic
allocation divides the portfolio.

\item[Stationary bootstrap] A bootstrap that preserves serial dependence
by drawing each observation from the sample or from its immediate
successor.

\item[Superior predictive ability] A bootstrap test that asks whether the
best member of a family beats a benchmark, with the family reweighted by
each member's own performance.

\item[Tilt] The adjustment the model makes to the composition of a sleeve.
Bounded, and never allowed to change a sleeve's total weight beyond a
stated scale.

\item[Wash sale] The rule that disallows a loss when the same or a
substantially identical security is bought within thirty days before or
after the sale. It applies in both directions and it carries the basis
and the holding period of the original position.

\end{description}
\clearpage

% ---- Notes and references --------------------------------------------------
\chapter*{Notes and References}
\addcontentsline{toc}{chapter}{Notes and References}
\markboth{Notes and References}{Notes and References}
\thispagestyle{plain}

\noindent\small
Every identifier this book cites was resolved against a registry before it
was used, and the registry's own title was compared with the title the
book intended to cite. The complete resolution table, with the registry
title for each source and the role it plays in the book, is Appendix~E,
and the ledger itself is published with the code.

Sources are referred to in the text by author and year. Where a claim
depends on a specific figure, the figure is given in the text and the
source is named in the same sentence.

\vspace{6pt}
\noindent\textbf{Behaviour.} Shiller (2015) for the account of market
exuberance the behavioural literature builds on; Odean (1998) for the
disposition effect and its magnitude; Barberis and Xiong (2009) for the modelled version; Barber
and Odean (2000) for the return of the most active retail traders; Kahneman
and Tversky (1979) for prospect theory; Benartzi and Thaler (1995) for
myopic loss aversion and the evaluation-period dependence of the equity
premium; Thaler (2008) for the sell-effect.

\vspace{6pt}
\noindent\textbf{Evaluation.} Bailey and L\'opez de Prado (2014) for the
Deflated Sharpe Ratio; Bailey, Borwein, L\'opez de Prado and Zhu (2016) for
the probability of backtest overfitting and the combinatorial construction
behind it; Lo (2002) for the sampling error of the Sharpe ratio; Newey and
West (1987) for the heteroskedasticity and autocorrelation consistent
covariance matrix; Andrews (1991) for the bandwidth rule; White (2000) for
the reality check; Hansen (2005) for the superior predictive ability test;
Abadie (2018) for the distinction between statistical and economic
significance; Sullivan, Timmermann and White (1999) for data snooping in
rule-based trading.

\vspace{6pt}
\noindent\textbf{Evidence.} Jensen, Kelly and Pedersen (2023) for the
replication evidence and the standard-error finding; McLean and Pontiff
(2016) for post-publication decay; Goyal, Welch and Zafirov (2024) for the
review of equity-premium predictors; Harvey, Liu and Zhu (2016) for the
multiple-testing hurdle; S\&P Dow Jones Indices for the annual SPIVA
survey of active managers.

\vspace{6pt}
\noindent\textbf{Portfolio construction.} Markowitz (1952); Kelly (1956)
for the criterion; Ledoit and Wolf (2004) for the shrinkage estimator; De
Miguel, Garlappi and Uppal (2009) for the failure of optimised portfolios
to beat equal weighting; L\'opez de Prado (2016) for hierarchical risk
parity; Campbell, Lo and MacKinlay (1997) for the econometrics of
returns; Hansen (2021) and Coqueret and Guida (2020) for the practice of
machine learning in asset management.

\vspace{6pt}
\noindent\textbf{Risk.} Moreira and Muir (2017) for volatility-managed
portfolios; Cederburg, O'Doherty, Wang and Yan (2020) for the replication
that finds the effect does not survive real-time implementation; Lo and
Zhang (2024) for the adaptive markets hypothesis; Samuelson (1969) for the
sequence of returns; Blanchett, Finke and Pfau (2013) for the safe
withdrawal rate under low bond yields.

\vspace{6pt}
\noindent\textbf{Factors and costs.} Fama and French (1993; 2015);
Asness, Moskowitz and Pedersen (2013); Novy-Marx (2013); Frazzini and
Pedersen (2014); Frazzini, Israel and Moskowitz (2012; 2018); Dong, Kang
and Peress (2025); Amihud (2002); Carhart (1997); Brown (1992); for the
live-versus-backtest boundary, the inception dates of the factor funds
themselves, which are listed in Appendix~C.

\vspace{6pt}
\noindent\textbf{Prior art.} Dhar and Sedoc (2025) for DBOT; Hwang and
Zohren (2025) for the signature-informed transformer; DeePM (2026) for
regime-robust deep learning; and the systematic review of deep
reinforcement learning for automated equity portfolio management (2026).

\vspace{6pt}
\noindent\textbf{Law and regulation.} Securities Exchange Act of 1940
Section 10(b) and Rule 10b-5; Investment Advisers Act of 1940 Sections
9(a)(2) and 9(a)(3); Regulation FD, 17~CFR 243.100--243.103; Financial
Industry Regulatory Authority Regulatory Notice 26-10 (2026); the
Securities and Exchange Commission's investor bulletin on exchange-traded
funds; and the Internal Revenue Code's wash-sale and capital-loss
provisions.

\vspace{6pt}
\noindent\small\textit{Every identifier above, with its registry title and
its role in the book, is listed in Appendix~E. The resolution ledger is
published with the source code, and the resolution script is part of it, so
that the claim can be re-checked rather than trusted.}
\clearpage

% ---- Appendix A -------------------------------------------------------------
\chapter{The Data and Its Provenance}
\label{app:data}

\thispagestyle{plain}
\addcontentsline{toc}{chapter}{Appendix A: The Data and Its Provenance}
\markboth{Appendix A}{Appendix A}

The book's data layer draws every price from a single source, caches the
raw responses with SHA-256 checksums, validates them, and builds the
panel from the cache. This appendix documents what was fetched, what was
checked, and what the checks found.

\section{The panel}

The monthly panel runs from \numPanelStart{} to \numPanelEnd{}: 407
months, 61 funds. The system's own record runs from \numRecordStart{} to
\numRecordEnd{}, which is \numRecordMonths{} months; the difference is the
warm-up the model needs before it will produce a forecast, plus the years
needed to establish the benchmark's own history.

Monthly returns are rebuilt from the raw close and the dividend cash paid
in that month, rather than read from the vendor's adjusted close, because
the tax layer needs a genuine transaction price and because the vendor
applies its dividend factor on a different convention. The two series
agree to within 0.004 per cent a month on average.

\section{The integrity checks}

Five checks are applied to every series, and a series that fails any of
them does not enter the panel: finite values; strictly positive prices;
strictly increasing dates with no duplicates; no gap longer than ten
calendar days; and no implausible single-day return.

All \numNFunds{} series in the universe pass all five checks. The raw
responses and their checksums are published with the code, so the
verification can be re-run against the cache without touching the
network.

\section{The cached series}

The table lists every raw response in the cache, with the first and last
bar and the number of bars. Sixty-three series were fetched; two of the
original sixty-five tickers --- a cryptocurrency product and a European
equity product --- do not exist at the vendor and were dropped rather than
substituted, which is why the universe has \numNFunds{} members.

\begingroup
\scriptsize
% ---------------------------------------------------------------------
%  Appendix A — generated from data/cache/manifest.json
% ---------------------------------------------------------------------
\begin{tabularx}{\textwidth}{@{}P{0.62in}P{0.92in}P{0.92in}Y@{}}
\toprule
Ticker & First bar & Last bar & Bars \\
\midrule
ACWI & 2008-03-28 & 2026-11-24 & 4696 \\
AVUV & 2019-09-26 & 2026-11-24 & 1801 \\
BIL & 2007-05-30 & 2026-11-24 & 4905 \\
BND & 2007-04-10 & 2026-11-24 & 4940 \\
BNDX & 2013-06-04 & 2026-11-24 & 3391 \\
DBC & 2006-02-06 & 2026-11-24 & 5234 \\
DBMF & 2019-05-08 & 2026-11-24 & 1899 \\
DEM & 2007-07-13 & 2026-11-24 & 4874 \\
EEM & 2003-04-14 & 2026-11-24 & 5943 \\
EFA & 2001-08-27 & 2026-11-24 & 6349 \\
EFV & 2005-08-05 & 2026-11-24 & 5360 \\
GLD & 2004-11-18 & 2026-11-24 & 5539 \\
HYG & 2007-04-11 & 2026-11-24 & 4939 \\
IAU & 2005-01-28 & 2026-11-24 & 5491 \\
IEF & 2002-07-30 & 2026-11-24 & 6121 \\
IVV & 2000-05-19 & 2026-11-24 & 6669 \\
IWM & 2000-05-26 & 2026-11-24 & 6664 \\
JNK & 2007-12-04 & 2026-11-24 & 4774 \\
LQD & 2002-07-30 & 2026-11-24 & 6121 \\
MDY & 1995-05-04 & 2026-11-24 & 7943 \\
MGK & 2007-12-27 & 2026-11-24 & 4758 \\
MGY & 2017-06-29 & 2026-11-24 & 2365 \\
MTUM & 2013-04-18 & 2026-11-24 & 3423 \\
MUB & 2007-09-10 & 2026-11-24 & 4834 \\
PDBC & 2014-11-07 & 2026-11-24 & 3029 \\
QQQ & 1999-03-10 & 2026-11-24 & 6972 \\
QUAL & 2013-07-18 & 2026-11-24 & 3360 \\
RINF & 2012-01-12 & 2026-11-24 & 3739 \\
RSP & 2003-05-01 & 2026-11-24 & 5931 \\
SFY & 2019-04-11 & 2026-11-24 & 1917 \\
SHV & 2007-01-11 & 2026-11-24 & 5000 \\
SHY & 2002-07-30 & 2026-11-24 & 6121 \\
SIZE & 2013-04-18 & 2026-11-24 & 3423 \\
SPDW & 2007-04-26 & 2026-11-24 & 4928 \\
SPHD & 2012-10-26 & 2026-11-24 & 3539 \\
SPY & 1993-01-29 & 2026-11-24 & 8514 \\
TIP & 2003-12-05 & 2026-11-24 & 5779 \\
TLT & 2002-07-30 & 2026-11-24 & 6121 \\
USMV & 2011-10-20 & 2026-11-24 & 3796 \\
USSG & 2019-03-08 & 2026-11-24 & 1941 \\
VCIT & 2009-11-23 & 2026-11-24 & 4277 \\
VEA & 2007-07-26 & 2026-11-24 & 4865 \\
VGLT & 2010-01-04 & 2026-11-24 & 4250 \\
VLUE & 2013-04-18 & 2026-11-24 & 3423 \\
VNQ & 2004-09-29 & 2026-11-24 & 5575 \\
VNQI & 2010-11-01 & 2026-11-24 & 4041 \\
VTI & 2001-06-15 & 2026-11-24 & 6399 \\
VTV & 2004-01-30 & 2026-11-24 & 5742 \\
VUG & 2004-01-30 & 2026-11-24 & 5742 \\
VWO & 2005-03-10 & 2026-11-24 & 5463 \\
XLB & 1998-12-22 & 2026-11-24 & 7024 \\
XLC & 2018-06-19 & 2026-11-24 & 2121 \\
XLE & 1998-12-22 & 2026-11-24 & 7024 \\
XLF & 1998-12-22 & 2026-11-24 & 7024 \\
XLI & 1998-12-22 & 2026-11-24 & 7024 \\
XLK & 1998-12-22 & 2026-11-24 & 7024 \\
XLP & 1998-12-22 & 2026-11-24 & 7024 \\
XLRE & 2015-10-08 & 2026-11-24 & 2799 \\
XLU & 1998-12-22 & 2026-11-24 & 7024 \\
XLV & 1998-12-22 & 2026-11-24 & 7024 \\
XLY & 1998-12-22 & 2026-11-24 & 7024 \\
\bottomrule
\end{tabularx}

\endgroup
\clearpage

% ---- Appendix B -------------------------------------------------------------
\chapter{The Estimators, From Their Objective Functions}
\label{app:estimators}

\thispagestyle{plain}
\addcontentsline{toc}{chapter}{Appendix B: The Estimators}
\markboth{Appendix B}{Appendix B}

The book's package depends on \code{numpy}, \code{pandas} and
\code{matplotlib}, and on nothing else. Every statistical estimator it
uses is written from its defining objective, including the two
distributions everything else is built on. This appendix lists what is
implemented and why.

\section{The distributions}

\code{patientquant/stats.py} implements the error function by its
definition as a numerical integral, the normal cumulative distribution
from the error function, the normal quantile function by Acklam's rational
approximation refined with two Newton--Halley steps, the incomplete beta
function by its continued-fraction expansion, the Student's \(t\)
cumulative distribution from the incomplete beta, and the Student's \(t\)
quantile function by bisection on the inverse.

\section{The model}

\code{patientquant/models.py} implements the normal inverse-gamma
conjugate update of Chapter 17 directly from the posterior formulas: the
posterior precision and mean, the scalar parameter updates, the
posterior variance of the mean, and the walk-forward training-set
construction. The online Bayesian model average of Chapter 19 is
implemented from its definition: a proper scoring rule on the standardised
forecast error, applied horizon-aware, accumulated with an exponential
discount, and combined by the law of total variance.

\section{The portfolio layer}

\code{patientquant/portfolio.py} implements the Ledoit--Wolf shrinkage
estimator from its definition --- the shrinkage intensity being the ratio
of the estimated squared deviation of the sample covariance from the
scaled identity to the estimated squared error of the sample covariance
itself --- the pairwise-complete covariance used when a fund has not been
listed for the whole estimation window, and the certainty-equivalent Kelly
solve.

\section{The evaluation machinery}

\code{patientquant/evaluation.py} implements the geometric and arithmetic
means, the compound growth rate, the maximum drawdown, the time to
recovery, the worst rolling twelve-month outcome, the Sortino and Calmar
ratios, Lo's scale factor \(\eta(q)\), the first-order autocorrelation of
a return series, and the Newey--West mean t-statistic with an
Andrews-style bandwidth.

\code{patientquant/selection.py} implements the expected maximum Sharpe
ratio in closed form and by Monte Carlo, the Deflated Sharpe Ratio, the
average pairwise correlation, the implied number of independent trials,
the combinatorial symmetric cross-validation that produces the
probability of backtest overfitting, the stationary bootstrap, and both
the White reality check and the Hansen superior predictive ability test.

\section{The tax engine}

\code{patientquant/costs.py} implements the lot accounting, the
first-in-first-out disposal, the holding-period classification, the
wash-sale disallowance in both directions with the basis adjustment and
the holding-period tack, the capital-loss offset against ordinary income
with its carryforward, and the annual settlement.

\begin{takeaway}[Why it is written this way]
A reader who wants to check the book's arithmetic can check it against
the definitions rather than against a library's implementation of them.
The estimators are short, the objective of each is stated in the module
docstring, and the tests in the repository exercise them against values
that can be computed independently.
\end{takeaway}
\clearpage

% ---- Appendix C -------------------------------------------------------------
\chapter{The Universe and Its Terms}
\label{app:universe}

\thispagestyle{plain}
\addcontentsline{toc}{chapter}{Appendix C: The Universe and Its Terms}
\markboth{Appendix C}{Appendix C}

Every fund in the universe, its sleeve, its expense ratio, the date it
listed, and the date from which it had a return on the panel. A fund
cannot be held before it has a price, so the last column is the boundary
of every backtest in the book.

The expense ratios are shown as percentages. They are not charged
separately in the simulation, because they are already inside the return
series; the mean across the universe is \numERMean\% a year, the cheapest
is \numERMin\% and the most expensive is \numERMax\%.

\begingroup
\scriptsize
% ---------------------------------------------------------------------
%  Appendix C — generated from build/series/fund_terms.csv
% ---------------------------------------------------------------------
\begin{tabularx}{\textwidth}{@{}P{0.46in}P{1.06in}P{0.62in}P{0.48in}P{0.50in}P{0.54in}Y@{}}
\toprule
Ticker & Name & Sleeve & Expense & Listed & Live from & First return \\
\midrule
BND & Vanguard Total Bond Market ETF & bond & 0.03\% & 2007-04-03 & 2007-04-30 & 236 months \\
BNDX & Vanguard Total International Bond ETF & bond & 0.07\% & 2013-05-31 & 2013-06-30 & 162 months \\
DBC & Invesco DB Commodity Index Tracking Fund & commodity & 0.85\% & 2006-02-03 & 2006-02-28 & 250 months \\
GLD & SPDR Gold Shares & commodity & 0.40\% & 2004-11-18 & 2004-11-30 & 265 months \\
IAU & iShares Gold Trust & commodity & 0.25\% & 2005-01-21 & 2005-01-31 & 263 months \\
PDBC & Invesco Optimum Yield Diversified Commodity Strategy & commodity & 0.59\% & 2014-11-07 & 2014-11-30 & 145 months \\
HYG & iShares iBoxx USD High Yield Corporate Bond ETF & credit & 0.49\% & 2007-04-04 & 2007-04-30 & 236 months \\
JNK & SPDR Bloomberg High Yield Bond ETF & credit & 0.40\% & 2007-11-28 & 2007-12-31 & 228 months \\
LQD & iShares iBoxx USD Investment Grade Corporate Bond ETF & credit & 0.14\% & 2002-07-22 & 2002-07-31 & 293 months \\
VCIT & Vanguard Intermediate-Term Corporate Bond ETF & credit & 0.03\% & 2009-11-19 & 2009-11-30 & 205 months \\
AVUV & Avantis US Value ETF & factor & 0.25\% & 2019-09-24 & 2019-09-30 & 87 months \\
DEM & Dynamic Equity ETF & factor & 0.63\% & 2007-07-13 & 2007-07-31 & 233 months \\
EFV & iShares MSCI EAFE Value Factor ETF & factor & 0.31\% & 2005-08-01 & 2005-08-31 & 256 months \\
MTUM & iShares MSCI USA Momentum Factor ETF & factor & 0.15\% & 2013-04-16 & 2013-04-30 & 164 months \\
QUAL & iShares MSCI USA Quality Factor ETF & factor & 0.15\% & 2013-07-16 & 2013-07-31 & 161 months \\
SFY & First Trust S\&P 500 Factor ETF & factor & 0.05\% & 2019-04-10 & 2019-04-30 & 92 months \\
SIZE & iShares MSCI USA Size Factor ETF & factor & 0.15\% & 2013-04-16 & 2013-04-30 & 164 months \\
SPHD & Invesco S\&P 500 High Dividend ETF & factor & 0.30\% & 2012-10-18 & 2012-10-31 & 170 months \\
USMV & iShares MSCI USA Min Vol Factor ETF & factor & 0.15\% & 2011-10-18 & 2011-10-31 & 182 months \\
USSG & iShares MSCI USA ESG Select ETF & factor & 0.09\% & 2019-03-06 & 2019-03-31 & 93 months \\
VLUE & iShares MSCI USA Value Factor ETF & factor & 0.15\% & 2013-04-16 & 2013-04-30 & 164 months \\
ACWI & iShares MSCI ACWI ETF & intl\_equity & 0.32\% & 2008-03-26 & 2008-03-31 & 225 months \\
EEM & iShares MSCI Emerging Markets ETF & intl\_equity & 0.72\% & 2003-04-07 & 2003-04-30 & 284 months \\
EFA & iShares MSCI EAFE ETF & intl\_equity & 0.32\% & 2001-08-14 & 2001-08-31 & 304 months \\
SPDW & SPDR MSCI ACWI IMI ETF & intl\_equity & 0.03\% & 2007-04-20 & 2007-04-30 & 236 months \\
VEA & Vanguard FTSE Developed Markets ETF & intl\_equity & 0.03\% & 2007-07-20 & 2007-07-31 & 233 months \\
VWO & Vanguard FTSE Emerging Markets ETF & intl\_equity & 0.06\% & 2005-03-04 & 2005-03-31 & 261 months \\
DBMF & Invesco Managed Futures ETF & managed\_futures & 0.85\% & 2019-05-07 & 2019-05-31 & 91 months \\
MUB & iShares National Muni Bond ETF & muni & 0.05\% & 2007-09-07 & 2007-09-30 & 231 months \\
RINF & Amplify PureAlpha Real Estate Income ETF & real\_estate & 0.30\% & 2012-01-10 & 2012-01-31 & 179 months \\
VNQ & Vanguard Real Estate ETF & real\_estate & 0.13\% & 2004-09-23 & 2004-09-30 & 267 months \\
VNQI & Vanguard International Real Estate ETF & real\_estate & 0.12\% & 2010-11-01 & 2010-11-30 & 193 months \\
XLB & Materials Select Sector SPDR & sector & 0.08\% & 1998-12-16 & 1998-12-31 & 336 months \\
XLC & Communication Services Select Sector SPDR & sector & 0.08\% & 2018-06-18 & 2018-06-30 & 102 months \\
XLE & Energy Select Sector SPDR & sector & 0.08\% & 1998-12-16 & 1998-12-31 & 336 months \\
XLF & Financial Select Sector SPDR & sector & 0.08\% & 1998-12-16 & 1998-12-31 & 336 months \\
XLI & Industrial Select Sector SPDR & sector & 0.08\% & 1998-12-16 & 1998-12-31 & 336 months \\
XLK & Technology Select Sector SPDR & sector & 0.08\% & 1998-12-16 & 1998-12-31 & 336 months \\
XLP & Consumer Staples Select Sector SPDR & sector & 0.08\% & 1998-12-16 & 1998-12-31 & 336 months \\
XLRE & Real Estate Select Sector SPDR & sector & 0.08\% & 2015-10-07 & 2015-10-31 & 134 months \\
XLU & Utilities Select Sector SPDR & sector & 0.08\% & 1998-12-16 & 1998-12-31 & 336 months \\
XLV & Health Care Select Sector SPDR & sector & 0.08\% & 1998-12-16 & 1998-12-31 & 336 months \\
XLY & Consumer Discretionary Select Sector SPDR & sector & 0.08\% & 1998-12-16 & 1998-12-31 & 336 months \\
TIP & iShares TIPS Bond ETF & tips & 0.18\% & 2003-12-04 & 2003-12-31 & 276 months \\
BIL & SPDR Bloomberg 1-3 Month T-Bill ETF & treasury & 0.14\% & 2007-05-25 & 2007-05-31 & 235 months \\
IEF & iShares 7-10 Year Treasury Bond ETF & treasury & 0.15\% & 2002-07-22 & 2002-07-31 & 293 months \\
SHV & iShares Short Treasury Bond ETF & treasury & 0.15\% & 2007-01-05 & 2007-01-31 & 239 months \\
SHY & iShares 1-3 Year Treasury Bond ETF & treasury & 0.15\% & 2002-07-22 & 2002-07-31 & 293 months \\
TLT & iShares 20+ Year Treasury Bond ETF & treasury & 0.15\% & 2002-07-22 & 2002-07-31 & 293 months \\
VGLT & Vanguard Long-Term Treasury ETF & treasury & 0.03\% & 2009-11-19 & 2010-01-31 & 203 months \\
IVV & iShares Core S\&P 500 ETF & us\_equity & 0.03\% & 2000-05-15 & 2000-05-31 & 319 months \\
IWM & iShares Russell 2000 ETF & us\_equity & 0.19\% & 2000-05-22 & 2000-05-31 & 319 months \\
MDY & SPDR S\&P MidCap 400 ETF Trust & us\_equity & 0.23\% & 1995-05-04 & 1995-05-31 & 379 months \\
MGK & VanEck Morningstar US Growth ETF & us\_equity & 0.05\% & 2007-12-17 & 2007-12-31 & 228 months \\
MGY & VanEck Morningstar US Value ETF & us\_equity & --- & --- & 2017-06-30 & 114 months \\
QQQ & Invesco QQQ Trust & us\_equity & 0.18\% & 1999-03-10 & 1999-03-31 & 333 months \\
RSP & Invesco S\&P 500 Equal Weight ETF & us\_equity & 0.20\% & 2003-04-24 & 2003-05-31 & 283 months \\
SPY & SPDR S\&P 500 ETF Trust & us\_equity & 0.09\% & 1993-01-22 & 1993-01-31 & 407 months \\
VTI & Vanguard Total Stock Market ETF & us\_equity & 0.03\% & 2001-05-24 & 2001-06-30 & 306 months \\
VTV & Vanguard Value ETF & us\_equity & 0.03\% & 2004-01-26 & 2004-01-31 & 275 months \\
VUG & Vanguard Growth ETF & us\_equity & 0.03\% & 2004-01-26 & 2004-01-31 & 275 months \\
\bottomrule
\end{tabularx}

\endgroup
\clearpage

% ---- Appendix D -------------------------------------------------------------
\chapter{The Verification Gates}
\label{app:gates}

\thispagestyle{plain}
\addcontentsline{toc}{chapter}{Appendix D: The Verification Gates}
\markboth{Appendix D}{Appendix D}

The book's build script runs a set of checks after the PDF is produced and
fails if any of them does. The checks exist because the failure modes they
catch are all invisible in the output: a stale figure, a hand-transcribed
number, a caption pointing at the wrong table, a page with overlapping
labels.

\begin{enumerate}
\item \textbf{The build completes.} The typesetting runs to a PDF with no
unresolved references, no missing figures and no errors in the log.
\item \textbf{Every figure exists.} Each figure named in the source is
present in the figures directory, and no figure in the directory is
unreferenced.
\item \textbf{Every figure is referenced once.} No figure is used twice
and none is orphaned.
\item \textbf{No label overlaps.} Every text label in every figure is
checked against every curve, marker, legend and other label in the same
figure, and the check fails if any pair overlaps.
\item \textbf{Every macro is defined.} Each \code{\textbackslash num}
macro used in the text is defined in the generated macro file, and every
macro name is checked to be a legal control word, since a name carrying a
digit would typeset the page and lose the number without any error.
\item \textbf{Every macro agrees with the data.} Each macro is
recomputed from \texttt{build/numbers.json} and compared with what the
text would print.
\item \textbf{No stray number.} Every percentage and every decimal in the
prose is checked against the canonical data, and any that cannot be
traced is reported for inspection.
\item \textbf{The citation keys resolve.} Every citation in the source
appears in the source ledger, and every ledger entry is cited.
\item \textbf{The ledger is clean.} Every identifier in the ledger
resolves against its registry with a matching title.
\item \textbf{The page count is right.} The page count written into the
colophon matches the PDF.
\item \textbf{The word count is right.} The word count written into the
colophon matches the source.
\item \textbf{No template residue survives.} The text contains no
unresolved template, no substituted field and no construction marker of any
kind.
\item \textbf{The copyright string is exact.} The string
\texttt{Copyright 2026 Tico Internet LLC} appears and nothing else does.
\item \textbf{No leaked secrets.} The repository contains no key, token
or password.
\item \textbf{The fonts are embedded.} The PDF's fonts are the house
fonts and are fully embedded.
\item \textbf{The PDF is well formed.} The file passes a structural
check.
\item \textbf{The data cache is intact.} Every cached series matches the
checksum in the manifest.
\item \textbf{The panel rebuilds.} The monthly panel rebuilt from the
cache matches the panel the numbers were computed from.
\item \textbf{The system reruns.} The whole system runs end to end from
the cache and reproduces the numbers.
\item \textbf{No look-ahead in the training set.} Every training row
satisfies the horizon constraint, checked directly on the data.
\item \textbf{The source is well formed.} Every heading argument is
closed and every chapter file's braces balance, because an unclosed
argument makes LaTeX swallow the text that follows it without saying so.
\end{enumerate}

\begin{takeaway}[Why the gates exist]
The book's own claim is that evaluation is the part that matters. A claim
about evaluation is worth nothing if the evaluation itself is not checked,
and each of these gates catches one way in which a book like this can be
subtly wrong without looking wrong.
\end{takeaway}
\clearpage

% ---- Appendix E -------------------------------------------------------------
\chapter{The Source Ledger}
\label{app:ledger}

\thispagestyle{plain}
\addcontentsline{toc}{chapter}{Appendix E: The Source Ledger}
\markboth{Appendix E}{Appendix E}

Every identifier this book cites, resolved against a registry, with the
registry's own title and the role the source plays. The resolution
script is published with the code, and the ledger can be rebuilt by
running it.

All fifty identifiers resolved with a matching registry title. Where an
identifier was initially wrong, the registry's own record for the work
was adopted in its place; where a work could not be resolved at all, it
was removed from the book and the claim it was to support was either
carried by the book's own computation or by a source that does resolve.

\begingroup
\scriptsize
% ---------------------------------------------------------------------
%  Appendix E — generated from research/source_ledger.md
% ---------------------------------------------------------------------
\begin{tabular}{@{}P{1.32in}P{2.00in}P{0.92in}P{0.36in}@{}}
\toprule
Identifier & Registry title & Role in the book & Match \\
\midrule
\code{10.3905/jpm.2014.40.5.094} & The Deflated Sharpe Ratio: Correcting for Selection Bias, Backtest Overfitting, and Non-Normality & Deflated Sharpe Ratio & yes \\
\code{10.21314/jcf.2016.322} & The probability of backtest overfitting & Probability of backtest overfitting & yes \\
\code{10.2469/faj.v58.n4.2453} & The Statistics of Sharpe Ratios & Sharpe ratio estimation error & yes \\
\code{10.1111/1468-0262.00152} & A Reality Check for Data Snooping & Reality check for data snooping & yes \\
\code{10.1198/073500105000000063} & A Test for Superior Predictive Ability & Superior predictive ability test & yes \\
\code{10.1093/rfs/hhv059} & … and the Cross-Section of Expected Returns & The multiple-testing hurdle & yes \\
\code{10.2139/ssrn.2249314} & …and the Cross-Section of Expected Returns & Review of the factor zoo & yes \\
\code{10.2307/1913610} & A Simple, Positive Semi-Definite, Heteroskedasticity and Autocorrelation Consistent Covariance Matrix & HAC standard errors for alpha & yes \\
\code{10.2307/2938229} & Heteroskedasticity and Autocorrelation Consistent Covariance Matrix Estimation & Bandwidth selection rule & yes \\
\code{10.3386/w24403} & Statistical Non-Significance in Empirical Economics & Economic vs statistical significance & yes \\
\code{10.1111/jofi.13249} & Is There a Replication Crisis in Finance? & The replication evidence & yes \\
\code{10.1111/jofi.12365} & Does Academic Research Destroy Stock Return Predictability? & Post-publication decay & yes \\
\code{10.1093/rfs/hhae044} & A Comprehensive 2022 Look at the Empirical Performance of Equity Premium Prediction & The predictability of the equity premium & yes \\
\code{10.1111/0022-1082.00163} & Data‐Snooping, Technical Trading Rule Performance, and the Bootstrap & Data snooping in rule-based trading & yes \\
\code{10.1111/j.1540-6261.1952.tb01525.x} & PORTFOLIO SELECTION* & Mean-variance optimisation & yes \\
\code{10.1002/j.1538-7305.1956.tb03809.x} & A New Interpretation of Information Rate & The Kelly criterion & yes \\
\code{10.3905/jpm.2004.110} & Honey, I Shrunk the Sample Covariance Matrix & Shrinkage of the covariance matrix & yes \\
\code{10.1093/rfs/hhm075} & Optimal Versus Naive Diversification: How Inefficient is the 1/ N Portfolio Strategy? & The case against optimisation & yes \\
\code{10.3905/jpm.2016.42.4.059} & Building Diversified Portfolios that Outperform Out of Sample & Hierarchical risk parity & yes \\
\code{10.1111/jofi.12513} & Volatility‐Managed Portfolios & The case for volatility scaling & yes \\
\code{10.1016/j.jfineco.2020.04.015} & On the performance of volatility-managed portfolios & The documented failure of volatility scaling & yes \\
\code{10.1093/oso/9780199681143.001.0001} & The Adaptive Markets Hypothesis & The adaptive markets hypothesis & yes \\
\code{10.2307/1926559} & Lifetime Portfolio Selection By Dynamic Stochastic Programming & The sequence of returns & yes \\
\code{10.3905/jwm.2013.16.2.055} & Low Bond Yields and Safe Portfolio Withdrawal Rates & Safe withdrawal rates & yes \\
\code{10.2478/fprj-2015-0002} & Just How Safe are ‘Safe Withdrawal Rates’ in Retirement? & The fragility of withdrawal rates & yes \\
\code{10.1093/acprof:oso/9780199959327.001.0001} & Asset Management & Factor investing & yes \\
\code{10.1111/0022-1082.00072} & Are Investors Reluctant to Realize Their Losses? & The disposition effect & yes \\
\code{10.1111/0022-1082.00226} & Trading Is Hazardous to Your Wealth: The Common Stock Investment Performance of Individual Investors & The cost of trading too much & yes \\
\code{10.2307/2118511} & Myopic Loss Aversion and the Equity Premium Puzzle & Myopic loss aversion & yes \\
\code{10.1111/j.1540-6261.2009.01448.x} & What Drives the Disposition Effect? An Analysis of a Long‐Standing Preference‐Based Explanation & The disposition effect, modelled & yes \\
\code{10.2307/1914185} & Prospect Theory: An Analysis of Decision under Risk & Prospect theory & yes \\
\code{10.1515/9781400865536} & Irrational Exuberance & Market exuberance & yes \\
\code{10.1016/0304-405X(93)90023-5} & Common risk factors in the returns on stocks and bonds & The factor model & yes \\
\code{10.1016/j.jfineco.2014.10.010} & A five-factor asset pricing model & The five-factor model & yes \\
\code{10.1111/jofi.12021} & Value and Momentum Everywhere & Global value and momentum premia & yes \\
\code{10.1016/j.jfineco.2013.01.003} & The other side of value: The gross profitability premium & The profitability premium & yes \\
\code{10.1016/j.jfineco.2013.10.005} & Betting against beta & The low-beta premium & yes \\
\code{10.2139/ssrn.3229719} & Trading Costs & Real trading costs of anomalies & yes \\
\code{10.2139/ssrn.2294498} & Trading Costs of Asset Pricing Anomalies & Trading costs of anomalies & yes \\
\code{10.1093/rfs/hhaf036} & Fast and Slow Arbitrage: The Predictive Power of (Persistent) Capital Flows for Factor Returns & Slow arbitrage & yes \\
\code{10.1016/s1386-4181(01)00024-6} & Illiquidity and stock returns: cross-section and time-series effects & Illiquidity as a priced factor & yes \\
\code{10.1111/j.1540-6261.1997.tb03808.x} & On Persistence in Mutual Fund Performance & Fund persistence and survivorship & yes \\
\code{10.1093/rfs/5.4.553} & Survivorship Bias in Performance Studies & Survivorship bias & yes \\
\code{10.48550/arXiv.2504.05639} & DBOT: Artificial Intelligence for Systematic Long-Term Investing & The closest prior work & yes \\
\code{10.48550/arXiv.2510.03129} & Signature-Informed Transformer for Asset Allocation & Decision-focused learning & yes \\
\code{10.48550/arXiv.2601.05975} & DeePM: Regime-Robust Deep Learning for Systematic Macro Portfolio Management & Regime-robust deep learning & yes \\
\code{10.1007/s10791-026-10336-1} & Systematic review of reinforcement learning for automated equity portfolio management from single agent to multi agent systems & Review of deep reinforcement learning & yes \\
\code{10.1515/9781400830213} & The Econometrics of Financial Markets & The econometrics of returns & yes \\
\code{10.1017/9781108883658} & Machine Learning for Asset Managers & Machine learning for asset managers & yes \\
\code{10.1201/9781003034858} & Machine Learning for Factor Investing & Machine learning for factors & yes \\
\bottomrule
\end{tabular}

\noindent\textbf{dsr}\quad Bailey, López de Prado \emph{The Journal of Portfolio Management} (2014). The Deflated Sharpe Ratio: Correcting for Selection Bias, Backtest Overfitting, and Non-Normality. \url{10.3905/jpm.2014.40.5.094}. Deflated Sharpe Ratio\par
\noindent\textbf{pbo}\quad Bailey, Borwein, López de Prado, Zhu \emph{The Journal of Computational Finance} (2016). The probability of backtest overfitting. \url{10.21314/jcf.2016.322}. Probability of backtest overfitting\par
\noindent\textbf{lo2002}\quad Lo \emph{Financial Analysts Journal} (2002). The Statistics of Sharpe Ratios. \url{10.2469/faj.v58.n4.2453}. Sharpe ratio estimation error\par
\noindent\textbf{white2000}\quad White \emph{Econometrica} (2000). A Reality Check for Data Snooping. \url{10.1111/1468-0262.00152}. Reality check for data snooping\par
\noindent\textbf{hansen2005}\quad Hansen \emph{Journal of Business \&amp; Economic Statistics} (2005). A Test for Superior Predictive Ability. \url{10.1198/073500105000000063}. Superior predictive ability test\par
\noindent\textbf{harvey2016}\quad Harvey, Liu, Zhu \emph{Review of Financial Studies} (2015). … and the Cross-Section of Expected Returns. \url{10.1093/rfs/hhv059}. The multiple-testing hurdle\par
\noindent\textbf{harvey2015}\quad Harvey, Liu, Zhu \emph{} (). …and the Cross-Section of Expected Returns. \url{10.2139/ssrn.2249314}. Review of the factor zoo\par
\noindent\textbf{newey1987}\quad Newey, West \emph{Econometrica} (1987). A Simple, Positive Semi-Definite, Heteroskedasticity and Autocorrelation Consistent Covariance Matrix. \url{10.2307/1913610}. HAC standard errors for alpha\par
\noindent\textbf{andrews1991}\quad Andrews \emph{Econometrica} (1991). Heteroskedasticity and Autocorrelation Consistent Covariance Matrix Estimation. \url{10.2307/2938229}. Bandwidth selection rule\par
\noindent\textbf{abadie2018}\quad Abadie \emph{} (). Statistical Non-Significance in Empirical Economics. \url{10.3386/w24403}. Economic vs statistical significance\par
\noindent\textbf{jkp2023}\quad JENSEN, KELLY, PEDERSEN \emph{The Journal of Finance} (2023). Is There a Replication Crisis in Finance?. \url{10.1111/jofi.13249}. The replication evidence\par
\noindent\textbf{mclean2016}\quad MCLEAN, PONTIFF \emph{The Journal of Finance} (2016). Does Academic Research Destroy Stock Return Predictability?. \url{10.1111/jofi.12365}. Post-publication decay\par
\noindent\textbf{goyal2024}\quad Goyal, Welch, Zafirov \emph{The Review of Financial Studies} (2024). A Comprehensive 2022 Look at the Empirical Performance of Equity Premium Prediction. \url{10.1093/rfs/hhae044}. The predictability of the equity premium\par
\noindent\textbf{sullivan1999}\quad Sullivan, Timmermann, White \emph{The Journal of Finance} (1999). Data-Snooping, Technical Trading Rule Performance, and the Bootstrap. \url{10.1111/0022-1082.00163}. Data snooping in rule-based trading\par
\noindent\textbf{markowitz1952}\quad Markowitz \emph{The Journal of Finance} (1952). Portfolio Selection. \url{10.1111/j.1540-6261.1952.tb01525.x}. Mean-variance optimisation\par
\noindent\textbf{kelly1956}\quad Kelly \emph{Bell System Technical Journal} (1956). A New Interpretation of Information Rate. \url{10.1002/j.1538-7305.1956.tb03809.x}. The Kelly criterion\par
\noindent\textbf{ledoit2004}\quad Ledoit, Wolf \emph{The Journal of Portfolio Management} (2004). Honey, I Shrunk the Sample Covariance Matrix. \url{10.3905/jpm.2004.110}. Shrinkage of the covariance matrix\par
\noindent\textbf{demiguel2009}\quad DeMiguel, Garlappi, Uppal \emph{Review of Financial Studies} (2007). Optimal Versus Naive Diversification: How Inefficient is the 1/N Portfolio Strategy?. \url{10.1093/rfs/hhm075}. The case against optimisation\par
\noindent\textbf{lopezdeprado2016}\quad López de Prado \emph{The Journal of Portfolio Management} (2016). Building Diversified Portfolios that Outperform Out of Sample. \url{10.3905/jpm.2016.42.4.059}. Hierarchical risk parity\par
\noindent\textbf{moreira2017}\quad MOREIRA, MUIR \emph{The Journal of Finance} (2017). Volatility-Managed Portfolios. \url{10.1111/jofi.12513}. The case for volatility scaling\par
\noindent\textbf{cederburg2020}\quad Cederburg, O’Doherty, Wang, Yan \emph{Journal of Financial Economics} (2020). On the Performance of Volatility-Managed Portfolios. \url{10.1016/j.jfineco.2020.04.015}. The documented failure of volatility scaling\par
\noindent\textbf{lozhang2024}\quad Lo, Zhang \emph{} (). The Adaptive Markets Hypothesis. \url{10.1093/oso/9780199681143.001.0001}. The adaptive markets hypothesis\par
\noindent\textbf{samuelson1969}\quad Samuelson \emph{The Review of Economics and Statistics} (1969). Lifetime Portfolio Selection by Dynamic Stochastic Programming. \url{10.2307/1926559}. The sequence of returns\par
\noindent\textbf{blanchett2013}\quad Blanchett, Finke, Pfau \emph{The Journal of Wealth Management} (2013). Low Bond Yields and Safe Portfolio Withdrawal Rates. \url{10.3905/jwm.2013.16.2.055}. Safe withdrawal rates\par
\noindent\textbf{drew2015}\quad Drew, Walk \emph{Financial Planning Research Journal} (2015). Just How Safe are 'Safe Withdrawal Rates' in Retirement?. \url{10.2478/fprj-2015-0002}. The fragility of withdrawal rates\par
\noindent\textbf{ang2014}\quad Ang \emph{} (). Asset Management. \url{10.1093/acprof:oso/9780199959327.001.0001}. Factor investing\par
\noindent\textbf{odean1998}\quad Odean \emph{The Journal of Finance} (1998). Are Investors Reluctant to Realize Their Losses?. \url{10.1111/0022-1082.00072}. The disposition effect\par
\noindent\textbf{barberodean2000}\quad Barber, Odean \emph{The Journal of Finance} (2000). Trading Is Hazardous to Your Wealth: The Common Stock Investment Performance of Individual Investors. \url{10.1111/0022-1082.00226}. The cost of trading too much\par
\noindent\textbf{benartzi1995}\quad Benartzi, Thaler \emph{The Quarterly Journal of Economics} (1995). Myopic Loss Aversion and the Equity Premium Puzzle. \url{10.2307/2118511}. Myopic loss aversion\par
\noindent\textbf{barberisxiong2009}\quad BARBERIS, XIONG \emph{The Journal of Finance} (2009). What Drives the Disposition Effect? An Analysis of a Long-Standing Preference-Based Explanation. \url{10.1111/j.1540-6261.2009.01448.x}. The disposition effect, modelled\par
\noindent\textbf{kahneman1979}\quad Kahneman, Tversky \emph{Econometrica} (1979). Prospect Theory: An Analysis of Decision under Risk. \url{10.2307/1914185}. Prospect theory\par
\noindent\textbf{shiller2015}\quad Shiller \emph{} (). Irrational Exuberance. \url{10.1515/9781400865536}. Market exuberance\par
\noindent\textbf{famafrench1993}\quad Fama, French \emph{Journal of Financial Economics} (1993). Common risk factors in the returns on stocks and bonds. \url{10.1016/0304-405x(93)90023-5}. The factor model\par
\noindent\textbf{famafrench2015}\quad Fama, French \emph{Journal of Financial Economics} (2015). A five-factor asset pricing model. \url{10.1016/j.jfineco.2014.10.010}. The five-factor model\par
\noindent\textbf{asness2013}\quad ASNESS, MOSKOWITZ, PEDERSEN \emph{The Journal of Finance} (2013). Value and Momentum Everywhere. \url{10.1111/jofi.12021}. Global value and momentum premia\par
\noindent\textbf{novymarx2013}\quad Novy-Marx \emph{Journal of Financial Economics} (2013). The other side of value: The gross profitability premium. \url{10.1016/j.jfineco.2013.01.003}. The profitability premium\par
\noindent\textbf{frazzinipedersen2014}\quad Frazzini, Pedersen \emph{Journal of Financial Economics} (2014). Betting against beta. \url{10.1016/j.jfineco.2013.10.005}. The low-beta premium\par
\noindent\textbf{frazzini2018}\quad Frazzini, Israel, Moskowitz \emph{} (). Trading Costs. \url{10.2139/ssrn.3229719}. Real trading costs of anomalies\par
\noindent\textbf{frazzini2012}\quad Frazzini, Israel, Moskowitz \emph{SSRN Electronic Journal} (2012). Trading Costs of Asset Pricing Anomalies. \url{10.2139/ssrn.2294498}. Trading costs of anomalies\par
\noindent\textbf{dong2025}\quad Dong, Kang, Peress \emph{The Review of Financial Studies} (2025). Fast and Slow Arbitrage: The Predictive Power of (Persistent) Capital Flows for Factor Returns. \url{10.1093/rfs/hhaf036}. Slow arbitrage\par
\noindent\textbf{amihud2002}\quad Amihud \emph{Journal of Financial Markets} (2002). Illiquidity and stock returns: cross-section and time-series effects. \url{10.1016/s1386-4181(01)00024-6}. Illiquidity as a priced factor\par
\noindent\textbf{carhart1997}\quad Carhart \emph{The Journal of Finance} (1997). On Persistence in Mutual Fund Performance. \url{10.1111/j.1540-6261.1997.tb03808.x}. Fund persistence and survivorship\par
\noindent\textbf{brown1992}\quad Brown, Goetzmann, Ibbotson, Ross \emph{Review of Financial Studies} (1992). Survivorship Bias in Performance Studies. \url{10.1093/rfs/5.4.553}. Survivorship bias\par
\noindent\textbf{dbot}\quad Dhar, Vasant, Sedoc, João \emph{arXiv} (2025). DBOT: Artificial Intelligence for Systematic Long-Term Investing. \url{10.48550/arXiv.2504.05639}. The closest prior work\par
\noindent\textbf{sit}\quad Hwang, Yoontae, Zohren, Stefan \emph{arXiv} (2025). Signature-Informed Transformer for Asset Allocation. \url{10.48550/arXiv.2510.03129}. Decision-focused learning\par
\noindent\textbf{deepm}\quad Wood, Kieran, Roberts, Stephen J., Zohren, Stefan \emph{arXiv} (2026). DeePM: Regime-Robust Deep Learning for Systematic Macro Portfolio Management. \url{10.48550/arXiv.2601.05975}. Regime-robust deep learning\par
\noindent\textbf{drlreview}\quad Rout, Acharya, Aithal, Sharma \emph{Discover Computing} (2026). Systematic review of reinforcement learning for automated equity portfolio management from single agent to multi agent systems. \url{10.1007/s10791-026-10336-1}. Review of deep reinforcement learning\par
\noindent\textbf{campbell1997}\quad Campbell, Lo, MacKinlay \emph{} (). The Econometrics of Financial Markets. \url{10.1515/9781400830213}. The econometrics of returns\par
\noindent\textbf{hansen2021}\quad López de Prado \emph{} (). Machine Learning for Asset Managers. \url{10.1017/9781108883658}. Machine learning for asset managers\par
\noindent\textbf{coqueret2020}\quad Coqueret, Guida \emph{} (). Machine Learning for Factor Investing. \url{10.1201/9781003034858}. Machine learning for factors\par

\endgroup
\clearpage

% ---- Index ------------------------------------------------------------------
\printindex

% ---- Acknowledgments --------------------------------------------------------
\chapter*{Acknowledgments}
\addcontentsline{toc}{chapter}{Acknowledgments}
\markboth{Acknowledgments}{Acknowledgments}
\thispagestyle{plain}

The data in this book comes from a single vendor's public chart endpoint,
and the institution behind it has made it possible for a book like this to
be written at all.

The decision to report a null result was not easy, and it belongs to the
owner of the project, who asked for the truth and then did not ask for it
to be softened.

\begin{keytakeaways}
\item The book is a work of analysis, and the analysis is published with
it, so that the reader can check the arithmetic rather than trust it.
\end{keytakeaways}
\clearpage

% ---- About the author --------------------------------------------------------
\chapter*{About the Author}
\addcontentsline{toc}{chapter}{About the Author}
\markboth{About the Author}{About the Author}
\thispagestyle{plain}

\noindent\textbf{Gerald Enrique Nelson Mc Kenzie} is the founder of Tico
Internet LLC and the author of this volume and its predecessors. He
writes about the quantitative parts of subjects that are usually treated as
narrative: near-Earth object characterisation and impact accountability,
lunar surface architecture, dark matter and the beginning of the universe,
and now the arithmetic of long-horizon investing.

The method is the same in each case. The argument is built from primary
sources, the quantitative claims carry their own arithmetic, the
instruments are published so that the arithmetic can be checked, and the
limits of the analysis are stated in the same voice as its conclusions.

\begin{keytakeaways}
\item Every book in this series stands on its own science and publishes
its own instruments.
\end{keytakeaways}
\clearpage

% ---- Colophon ---------------------------------------------------------------
\chapter*{Colophon}
\addcontentsline{toc}{chapter}{Colophon}
\markboth{Colophon}{Colophon}
\thispagestyle{plain}

\noindent\small
\setmainfont{Outfit}
\noindent The Patient Quant is set in Outfit, with DejaVu Sans Mono for
code and identifiers and Fira Math for mathematics, and typeset with
XeLaTeX in a 6-by-9-inch trim. The figures are computed with Python and
matplotlib and embedded as vector graphics.

\noindent The body text runs to \bookWords{} words across \bookPages{}
numbered pages, in seven parts and thirty-eight chapters, with eighteen
figures, five appendices and an index.

\noindent Every number the book states about its own results is written by
\texttt{build/make\_numbers\_tex.py} from the system's output, and every
figure is drawn by \texttt{build/make\_figures.py} from the same files.
The verification gates of Appendix D are run by \texttt{build/verify.py}
after the typesetting, and the build fails if any of them does.

\vspace{8pt}
\noindent\textbf{DOI}\quad \href{https://doi.org/\bookDOI}{\texttt{\bookDOI}}\par

\vspace{8pt}
\noindent Copyright \textcopyright\ 2026 Tico Internet LLC.
\clearpage
